\chapter{The H12 and X30Controller}
\label{ch:u-handheld}

\begin{center}
\includegraphics[width=\linewidth]{h12}
\end{center}

\section{Sticks, keys and modes}

\begin{center}
\begin{tabularx}{\linewidth}{@{}>{\bfseries}p{0.3\linewidth}X@{}}
  \toprule
  Control & What it does \\
  \midrule
  Left stick & forward and back; left and right: step sideways \\
  Right stick & turn left and right \\
  \key{A} & AUTO or REMOTE \\
  \key{B} & STOP: the software e-stop \\
  \key{C} / \key{D} & slower / faster (three speed levels) \\
  Switch \key{E} & the speed level directly \\
  \bottomrule
\end{tabularx}
\end{center}

\begin{itemize}
  \item \textbf{REMOTE}: the sticks drive the robot. Nothing else does.
  \item \textbf{AUTO}: the robot follows its route; the sticks do nothing. The screen says
    \ui{AUTO mode: switch to REMOTE to drive} when you try.
\end{itemize}

The robot listens to one handheld at a time. When the H12 falls silent (out of range, app
closed), the robot stops its stick motion within a fraction of a second.

\section{The drive screen}

\begin{itemize}
  \item \textbf{Top bar}: battery, signal, mode, localization (\ui{Start pose}, \ui{Tag},
    \ui{No loc}), the protective field (\ui{Path clear}, \ui{Slowed}, \ui{Blocked},
    \ui{No scan}), health (\ui{Needs attention}) and the trip.
  \item \textbf{Large view}: the inspection camera, visible or thermal (the camera button
    switches); swap it with the map, which shows the route, its stop points and the robot.
  \item \textbf{STOP}: the red button, same as \key{B}.
  \item \textbf{Actions}: \ui{Record}, \ui{Point}, \ui{Finish}, \ui{Tags}, \ui{Localize},
    \ui{Routes}, \ui{Start}, \ui{Pause}, \ui{Resume}, \ui{End}, \ui{Inspect}, the recorder, and
    the robot's postures (stand up, sit down).
\end{itemize}

\section{STOP and release}

STOP (key \key{B} or the red button) holds the robot still: \ui{E-stop on: the robot is held
still}. A patrol pauses. To release, \textbf{hold} STOP for a second (\ui{Hold to release}).
After a release the robot moves only once the sticks were let go: a stick held during the
release never moves it.

\section{On a phone or tablet}

The same app (version 1.1 or later) runs on any Android phone or tablet. Without the H12 it shows
touch sticks on the screen: the left one drives and steps sideways, the right one turns; let go
and the robot stops. \menupath{Settings › Sticks} chooses who drives:

\begin{center}
\begin{tabularx}{\linewidth}{@{}>{\bfseries}p{0.2\linewidth}X@{}}
  \toprule
  Choice & Drives \\
  \midrule
  Auto & the H12's sticks, else a gamepad, else the touch sticks \\
  H12 & only the H12's sticks \\
  Gamepad & only a gamepad (Bluetooth or USB) \\
  Touch & only the touch sticks \\
  \bottomrule
\end{tabularx}
\end{center}
